\chapter*{Declaração de Uso de Inteligência Artificial Generativa}

Declaro que ferramentas de Inteligência Artificial Generativa foram utilizadas no desenvolvimento deste trabalho. Abaixo, detalha-se o uso específico das tecnologias empregadas.

\vspace{1em}

\noindent\textbf{Ferramenta/Versão:} Claude Sonnet 4.8.

\noindent\textbf{Etapa da pesquisa:} desenvolvimento da aplicação e apoio à implementação.

\noindent\textbf{Finalidade do uso:} apoiar a estruturação de componentes da solução, auxiliar na geração e revisão de trechos de código, depuração, organização de fluxos de implementação e análise técnica de decisões de desenvolvimento.

\vspace{0.8em}

\noindent\textbf{Ferramenta/Versão:} ChatGPT 5.5.

\noindent\textbf{Etapa da pesquisa:} revisão textual, organização do documento e adaptação para \LaTeX{}.

\noindent\textbf{Finalidade do uso:} auxiliar na melhoria da clareza, fluidez e organização do texto acadêmico, bem como na conversão e adequação de seções ao formato utilizado no Overleaf com a classe \textit{infufrgs}. Adicionalmente, foi utilizado na elaboração das personas de pacientes simulados (Ana, Carlos e Maria), apoiando a definição de características comunicacionais coerentes com os atributos de cada perfil.

\vspace{0.8em}

\noindent\textbf{Ferramenta/Versão:} Gemini 2.5 Flash-Lite e Gemini 2.5 Flash. 

\noindent\textbf{Etapa da pesquisa:} execução experimental da solução SimplyHealth. 

\noindent\textbf{Finalidade do uso:} utilização como modelo de linguagem integrado à solução desenvolvida, responsável pela geração de versões simplificadas dos textos de saúde e por etapas de avaliação automatizada relacionadas aos guardrails.
\vspace{0.8em}

\noindent\textbf{Nota de responsabilidade:} o autor revisou criticamente todo o conteúdo gerado ou modificado pelas ferramentas mencionadas, assumindo responsabilidade exclusiva e integral pela originalidade, integridade e veracidade das informações apresentadas no trabalho final.
